% What's the problem?
% Why is it a problem? Research gap left by other approaches?
% Why is it important? Why care?
% What's the approach? How to solve the problem?
% What's the findings? How was it evaluated, what are the results, limitations, 
% what remains to be done?

% XXX Summary
\emph{Summary:}
\dots

% Motivation and intended learning outcomes
\emph{Avsedda lärandemål:}
Efter att ha genomfört detta kapitel ska du kunna:
\begin{itemize}
  \item Förklara hur ett grafiskt program drivs av en händelseslinga och
    vad en \foreignlanguage{english}{callback}-funktion är. % Lo08
  \item Konstruera ett enkelt grafiskt gränssnitt med
    \mintinline{python}{tkinter} med fönster, etiketter, textfält och
    knappar. % Lo05, Lo08
  \item Koppla händelser från användaren (knapptryck, tangenter) och från
    klockan till funktioner och metoder som uppdaterar
    gränssnittet. % Lo08
  \item Strukturera ett grafiskt program med klasser, till exempel genom
    att ärva från \mintinline{python}{tkinter}-klasser, så att gränssnitt
    och programlogik hålls isär. % Lo03, Lo09
  \item Hitta och använda dokumentationen för ett programbibliotek för
    att lösa ett problem du inte sett förut. % Lo06
\end{itemize}

% XXX Prerequisites
\emph{Prerequisites:}
\dots

% XXX Reading material
\emph{Reading:}
\dots
